\documentclass[journal]{IEEEtran}

\usepackage{amsmath,amssymb,amsfonts}
\usepackage{booktabs}
\usepackage{graphicx}
\usepackage{cite}
\usepackage{multirow}
\usepackage{url}

\providecommand{\TODO}[1]{\textbf{[#1]}}

\begin{document}

\title{Recognize, Don't Generate: Language Pretraining as a Temporal-Structure
Prior in LLM Time-Series Forecasting}

\author{Anonymous% \IEEEauthorblockN{}
% \IEEEauthorblockA{}
}

\markboth{ICASSP 2027 Submission}{}

\maketitle

\begin{abstract}
Whether language pretraining contributes anything to LLM-based time-series
forecasting beyond architecture and tokenizer capacity remains contested.  We
study this question under a strict control: the same Qwen3-8B architecture,
tokenizer, data windows, and LoRA protocol, with only the initialization
differing (pretrained versus random).  Fine-tuning the pretrained model reduces
forecast MSE by $2.2\times$ on synthetic sine, $\sim 25\times$ on ETTm1, and
$\sim 6\times$ on ETTh1 (relative reductions of 125\%, 2380\%, and 495\%),
directly countering the claim that language pretraining is useless for
forecasting.  We then ask \emph{why} it helps.  A linear probe on the
\emph{frozen} LLM's hidden state classifies temporal dynamics (trend / periodic
/ local / mixture / regime) with 96.3\% accuracy, while on the \emph{same
windows} frozen numeric generation only produces 26.7\% complete forecasts with
$5.5\times$ the oracle error.  Language pretraining thus acts as a
temporal-structure prior: it helps the LLM \emph{recognize} dynamics more than
it makes it a better numeric \emph{generator}.  We exploit this with a frozen-LLM
expert router that selects among lightweight statistical experts: on
\emph{unseen} dynamics it matches the oracle $67.8\%$ of the time and reduces
forecast error by $\sim 42\%$ relative to a hand-feature router, while a
random-initialized same-architecture control transfers far worse.
\end{abstract}

\section{Introduction}
\label{sec:intro}
LLM-based time-series forecasting (LLM4TS) is an active area
\cite{timellm,autotimes,prdts}.  A central open question is \emph{why} language
pretraining should help predict numbers at all.  A prominent skeptical study
\cite{tan2024useful} found that removing or randomly initializing the LLM in
several popular LLM4TS methods does not hurt forecasting, concluding that the
gains come from architecture (patching, attention) rather than language
pretraining.  Recent controlled comparisons report the opposite for
parameter-efficient fine-tuning of modern LLMs \cite{randominit,crossmodalm}.

This debate is confounded by uncontrolled factors: methods freeze or tune
different components, use different tokenizers and prompts, and compare against
baselines of different capacity.  We contribute a controlled measurement and a
mechanistic hypothesis.  Under a strict control --- identical architecture,
tokenizer, windows, and LoRA protocol, with initialization as the only factor
--- we find that language pretraining helps forecast error \emph{and} that it
helps structure \emph{recognition} far more than numeric \emph{generation}.
Specifically:

\begin{itemize}
\item \textbf{Controlled pretraining ablation} (Sec.~\ref{sec:e1},
  Table~\ref{tab:forecast}): pretrained Qwen3-8B beats a random initialization
  of the same architecture by $2.2\times$--$25\times$ in MSE across
  sine/ETTm1/ETTh1, with perfect parse coverage.
\item \textbf{Structure recognition without generation} (Sec.~\ref{sec:e2}):
  a linear probe on the \emph{frozen} pretrained LLM classifies temporal
  dynamics at 96.3\% accuracy; a random same-architecture control reaches only
  90.6\% (and below 80\% on the hardest class, local).
\item \textbf{The asymmetry} (Sec.~\ref{sec:e3}): on the same windows,
  recognition is 96.7\% accurate while frozen numeric generation parses only
  26.7\% of windows and errs $5.5\times$ above the oracle.
\item \textbf{Zero-shot frozen-LLM router} (Sec.~\ref{sec:e4}): the probe
  doubles as an expert router on \emph{unseen} dynamics, matching the oracle
  $67.8\%$ of the time versus $12.8\%$ for a hand-feature router and $28.9\%$
  for the random control.
\end{itemize}

\section{Controlled Setup}
\label{sec:setup}
\textbf{Model.}  We use the local Wave-MoE Qwen3-8B checkpoint
(Qwen3TS-8B); only the language backbone is exercised.  \textbf{Text format.}
Each normalized history window is rendered as
``Forecast the next values of this normalized time series. History:
\texttt{+0.12 -0.04 ...} Forecast:'' with values clamped to $[-9.99,9.99]$
and formatted with two decimals.  \textbf{Training.}  LoRA (rank 8) on
$q/k/v/o$ projections; teacher forcing with the prompt tokens masked from the
LM loss; free-running greedy evaluation; context 64, horizon 16; 256 train / 64
test windows; 5 epochs; AdamW; train-only normalization.  \textbf{Control.}
The random initialization instantiates the identical architecture with random
weights and applies the same LoRA adapters, windows, and seed.
\textbf{Data.}  Synthetic sine, and the $OT$ channel of ETTm1 and ETTh1.

\textbf{Probe.}  For a frozen (no-tuning) model, we take the last-layer hidden
state of the final prompt token and train a linear (softmax) probe.  The
feature floor uses nine hand-crafted statistics (linear slope, fit $R^2$,
volatility, lag-1 autocorrelation, dominant period and strength, spectral
entropy, spectral peak concentration, range ratio).

\textbf{Dynamics data.}  Five synthetic dynamics kinds: \emph{trend} (linear
drift), \emph{periodic} (sinusoid), \emph{local} (AR(1)), \emph{mixture}
(trend + sinusoid), and \emph{regime} (trend switching to periodic mid-window).
Experts are lightweight and fit per window: \emph{trend} (linear extrapolation),
\emph{periodic} (autocorrelation period detection + seasonal-naive repeat), and
\emph{local} (stable damped AR($5$)).

\section{Experiments}

\subsection{Language pretraining helps: controlled ablation}
\label{sec:e1}
Table~\ref{tab:forecast} aggregates three seeds.  Pretrained initialization
reduces MSE by $2.2\times$ (sine), $25\times$ (ETTm1), and $6\times$ (ETTh1)
relative to random initialization of the same architecture, with parse
coverage $1.0$ and much lower variance; spectral amplitude error is also
halved to a third.  This directly contradicts \cite{tan2024useful} in the
parameter-efficient fine-tuning regime of a modern LLM.

\begin{table}[t]
\caption{Controlled forecast: pretrained vs.\ random initialization (mean
$\pm$ std over seeds $\{7,17,27\}$).  Rel.\ MSE = relative reduction.}
\label{tab:forecast}
\centering
\begin{tabular}{lcccccc}
\toprule
Dataset & Init & MSE & MAE & Parse & Spec.\ amp MAE & Rel.\\
\midrule
sine    & pretrained & 0.0094$\pm$0.0006 & 0.077 & 1.000 & 0.238 & \\
sine    & random     & 0.0212$\pm$0.0011 & 0.111 & 1.000 & 0.342 & 125\%\\
ETTm1   & pretrained & 0.0099$\pm$0.0017 & 0.071 & 1.000 & 0.215 & \\
ETTm1   & random     & 0.2453$\pm$0.1366 & 0.339 & 1.000 & 0.819 & 2380\%\\
ETTh1   & pretrained & 0.0487$\pm$0.0097 & 0.160 & 1.000 & 0.449 & \\
ETTh1   & random     & 0.2896$\pm$0.2059 & 0.366 & 0.990 & 0.876 & 495\%\\
\bottomrule
\end{tabular}
\end{table}

\subsection{Frozen structure recognition}
\label{sec:e2}
Table~\ref{tab:probe} reports 5-way dynamics classification on held-out
windows.  The frozen pretrained probe reaches 96.3\%, close to the hand-feature
floor (99.0\%) and above the random same-architecture control (90.6\%, weakest
on \emph{local}).  Clean synthetic dynamics are easy for every method; the
advantage of the pretrained representation appears under
distribution shift (Sec.~\ref{sec:e4}).

\begin{table}[t]
\caption{Dynamics recognition accuracy (5 classes, frozen).}
\label{tab:probe}
\centering
\begin{tabular}{lccccc}
\toprule
Method & Overall & Trend & Periodic & Local & Mixture\\
\midrule
Hand features + logistic & 0.988 & 1.00 & 0.97 & 0.98 & 1.00\\
Frozen pretrained Qwen & 0.963 & 1.00 & 0.94 & 0.94 & 0.94\\
Frozen random Qwen & 0.906 & 1.00 & 0.96 & 0.78 & 0.88\\
\bottomrule
\end{tabular}
\end{table}

\subsection{Recognize versus generate, on the same windows}
\label{sec:e3}
On 40 windows stratified across the five dynamics, the frozen pretrained probe
recognizes the dynamics kind with 96.7\% accuracy.  On the very same windows,
frozen greedy numeric generation parses only 26.7\% of the 16-step forecasts,
and the complete ones have MSE 1.85 --- $5.5\times$ the oracle expert error
(0.33).  Language pretraining therefore confers structure \emph{recognition};
numeric \emph{generation} is the fragile mode.

\subsection{Zero-shot frozen-LLM expert router}
\label{sec:e4}
We train the router on the three clean dynamics kinds and apply it to the two
\emph{unseen} kinds (mixture, regime).  Table~\ref{tab:router} compares the
hand-feature router, the frozen pretrained probe router, and the frozen random
probe router against the oracle, best-single-expert, and uniform baselines.
The pretrained probe matches the oracle 67.8\% of the time and reduces
downstream MSE by 42\% relative to the feature router (0.739 vs.\ 1.278); soft
routing improves this to 0.723.  The random control transfers far worse
(28.9\%, MSE 1.08), isolating language pretraining as the source of transfer.
The feature router collapses on unseen dynamics (12.8\%).  We note honestly
that on this hard set the best single expert (trend) remains competitive
(0.500) because trend dominates the oracle; routing's clear advantage is in
\emph{recognition transfer}, not in uniformly beating every expert.

\begin{table}[t]
\caption{Zero-shot router on unseen dynamics (mixture/regime).}
\label{tab:router}
\centering
\begin{tabular}{lccc}
\toprule
Method & Oracle-acc & MSE & Soft MSE\\
\midrule
Oracle (upper bound) & --- & 0.449 & ---\\
Best single (trend) & --- & 0.500 & ---\\
Uniform ensemble & --- & 1.287 & ---\\
Feature router & 0.128 & 1.278 & 1.235\\
Frozen random probe & 0.289 & 1.079 & 1.056\\
Frozen pretrained probe & 0.678 & 0.739 & 0.723\\
\bottomrule
\end{tabular}
\end{table}

\subsection{Real-data validation: zero-shot routing on ETTm1}
\label{sec:e5}
We apply the synthetic-trained routers to 300 test windows of the ETTm1 $OT$
channel without any ETTm1 training (Table~\ref{tab:ettm1}).  The hand-feature
router collapses: it routes 271 of 300 windows to the \emph{local} expert
(which has MSE 0.72 there), giving 6.0\% oracle accuracy and MSE 0.657, worse
than the uniform ensemble.  Both LLM-based routers --- pretrained and random
same-architecture --- avoid this trap entirely and select \emph{periodic} for
most windows (the best univariate expert, MSE 0.014).  This shows the
architecture and numeric-text tokenizer already encode enough structure to
recognize ETTm1's periodicity dominance.  However, because ETTm1's dynamics
are homogeneous, a degenerate router that always selects \emph{periodic}
already achieves oracle accuracy 67.7\% and the best-single MSE, and the random
probe collapses to exactly this degenerate policy.  The pretrained probe is
more discriminative (it also identifies some trend windows) but does not
improve downstream MSE on this homogeneous data.  The value of language
pretraining for routing therefore appears where dynamics are diverse and
\emph{unseen} (Sec.~\ref{sec:e4}), not on a single-homogeneous real series.
The same pattern reproduces on ETTh1 ($OT$): the feature router again collapses
into \emph{local} (MSE 0.748, 5\% oracle accuracy), both LLM-based routers
avoid it (MSE $\approx$ 0.040), and the random probe again degenerates to
always-$\emph{periodic}$.

\begin{table}[t]
\caption{Zero-shot expert routing on ETTm1 (300 test windows, no adaptation).}
\label{tab:ettm1}
\centering
\begin{tabular}{lccc}
\toprule
Method & Oracle-acc & MSE & Route (T/P/L)\\
\midrule
Oracle (4 experts) & --- & 0.010 & ---\\
Best single (periodic) & --- & 0.014 & ---\\
Always periodic (degenerate) & 0.677 & 0.014 & 0/300/0\\
Uniform ensemble & --- & 0.255 & ---\\
Feature router & 0.060 & 0.657 & 11/18/271\\
Frozen random probe & 0.677 & 0.014 & 0/300/0\\
Frozen pretrained probe & 0.517 & 0.020 & 82/218/0\\
\bottomrule
\end{tabular}
\end{table}

\section{Related Work}
\label{sec:related}
LLM4TS methods reprogram or fine-tune LLMs for forecasting
\cite{timellm,autotimes,prdts}; the skepticism line argues gains come from
architecture rather than language \cite{tan2024useful}.  Same-architecture
initialization studies \cite{randominit,crossmodalm} report pretraining
benefits under fine-tuning.  LLM-as-router and MoE hybrids for time series
\cite{lego,fmllm,dmoe} route tokens or tasks through experts; we instead route
\emph{temporal dynamics} with a \emph{frozen} LLM and measure zero-shot
transfer to unseen dynamics.  Our spectral diagnostics are natural for ICASSP.

\section{Conclusion}
\label{sec:conclusion}
Under a strict controlled comparison, language pretraining is a
temporal-structure prior in LLM time-series forecasting: it drives large
forecast-error gains under LoRA fine-tuning, and a \emph{frozen} pretrained
LLM recognizes temporal dynamics and routes experts zero-shot far better than
hand features or a random same-architecture control --- while numeric
generation remains the fragile mode.  Limitations: small window counts and
synthetic probe data (two seeds); clean stationary dynamics are easy for cheap
features, so the LLM's advantage is specifically under distribution shift.
Future work extends dynamic (patch-level) routing and real multivariate data.

\bibliographystyle{IEEEtran}
\bibliography{refs}

\end{document}
